\documentclass[10pt,a4paper,twocolumn]{article}

\usepackage[utf8]{inputenc}
\usepackage[indonesian]{babel}
\usepackage{amsmath,amssymb,amsfonts}
\usepackage{geometry}
\usepackage{tcolorbox}
\usepackage{booktabs}
\usepackage{enumitem}
\usepackage{hyperref}
\usepackage{xcolor}

% Margin ultra-efisien untuk tepat 1 lembar bolak-balik (2 halaman A4)
\geometry{
    top=0.8cm,
    bottom=0.8cm,
    left=0.85cm,
    right=0.85cm
}
\setlength{\columnsep}{4.2mm}
\pagestyle{empty}

\tcbuselibrary{skins}

% Warna konsisten & kontras tinggi untuk dicetak/disalin tangan
\definecolor{primaryblue}{RGB}{20, 60, 110}
\definecolor{secondaryblue}{RGB}{41, 128, 185}
\definecolor{boxbg}{RGB}{246, 249, 253}
\definecolor{darkgreen}{RGB}{24, 120, 64}
\definecolor{greenbg}{RGB}{243, 250, 245}
\definecolor{darkred}{RGB}{175, 30, 30}
\definecolor{redbg}{RGB}{253, 245, 245}
\definecolor{darkamber}{RGB}{185, 100, 10}
\definecolor{amberbg}{RGB}{254, 250, 242}

\newcommand{\angl}[1]{\overline{#1|}}

\newtcolorbox{cbox}[2][primaryblue]{
    enhanced,
    colback=boxbg,
    colframe=#1,
    fonttitle=\bfseries\sffamily\footnotesize,
    title={#2},
    boxrule=0.6pt,
    titlerule=0pt,
    coltitle=white,
    colbacktitle=#1,
    arc=1mm,
    left=1.8mm,
    right=1.8mm,
    top=1.1mm,
    bottom=1.1mm,
    before skip=1.2mm,
    after skip=1.5mm
}

\newtcolorbox{gbox}[1]{
    enhanced,
    colback=greenbg,
    colframe=darkgreen,
    fonttitle=\bfseries\sffamily\footnotesize,
    title={#1},
    boxrule=0.6pt,
    titlerule=0pt,
    coltitle=white,
    colbacktitle=darkgreen,
    arc=1mm,
    left=1.8mm,
    right=1.8mm,
    top=1.1mm,
    bottom=1.1mm,
    before skip=1.2mm,
    after skip=1.5mm
}

\newtcolorbox{rbox}[1]{
    enhanced,
    colback=redbg,
    colframe=darkred,
    fonttitle=\bfseries\sffamily\footnotesize,
    title={#1},
    boxrule=0.6pt,
    titlerule=0pt,
    coltitle=white,
    colbacktitle=darkred,
    arc=1mm,
    left=1.8mm,
    right=1.8mm,
    top=1.1mm,
    bottom=1.1mm,
    before skip=1.2mm,
    after skip=1.5mm
}

\newtcolorbox{abox}[1]{
    enhanced,
    colback=amberbg,
    colframe=darkamber,
    fonttitle=\bfseries\sffamily\footnotesize,
    title={#1},
    boxrule=0.6pt,
    titlerule=0pt,
    coltitle=white,
    colbacktitle=darkamber,
    arc=1mm,
    left=1.8mm,
    right=1.8mm,
    top=1.1mm,
    bottom=1.1mm,
    before skip=1.2mm,
    after skip=1.5mm
}

\setlength{\parindent}{0pt}
\setlength{\parskip}{1.2pt}

\hypersetup{
    colorlinks=true,
    linkcolor=primaryblue,
    pdftitle={AK2161 - Cheatsheet 1 Lembar Bolak-Balik Kuis Obligasi}
}

\begin{document}
\footnotesize

% ==============================================================================
% HALAMAN 1: SISI DEPAN (TIER 1 — RUMUS INTI & PALING SERING KELUAR)
% ==============================================================================
\twocolumn[{
\begin{tcolorbox}[
    enhanced,
    colback=primaryblue,
    colframe=primaryblue,
    arc=1.2mm,
    boxrule=0pt,
    left=3mm, right=3mm, top=1.6mm, bottom=1.6mm,
    after skip=2mm
]
\begin{center}
    {\color{white}\normalsize\textbf{\textsf{CHEATSHEET 1 LEMBAR BOLAK-BALIK --- SISI DEPAN (TIER 1: WAJIB \& SERING KELUAR)}}}\\[0.15em]
    {\color{amberbg}\scriptsize\textsf{AK2161 Teori Suku Bunga (Bab 12--14) $\bullet$ Urutan Prioritas Berdasarkan Arsip Kuis \& UAS ITB (2013--2026)}}
\end{center}
\end{tcolorbox}
}]

% KOTAK 1
\begin{cbox}{1. Kamus Simbol, Faktor Anuitas \& Konversi Frekuensi}
\textbf{Semua $r, i, n$ WAJIB dalam satuan periode kupon ($m$ kali/tahun)!}\\[0.8pt]
$\bullet$ $F$ = par value;\quad $C$ = redemption value;\quad $P = B_0$ = harga awal\\
$\bullet$ $r = \frac{\text{kupon nominal}}{m}$;\quad Kupon $= Fr$;\quad $g = \dfrac{Fr}{C}$ ($g = r$ jika $C = F$)\\
$\bullet$ $i$ = yield per periode kupon;\quad $v = (1+i)^{-1}$;\quad $K = Cv^n$;\quad $G = \dfrac{Fr}{i}$\\
$\bullet$ $a_{\angl{n}i} = \dfrac{1 - (1+i)^{-n}}{i}$, \qquad $s_{\angl{n}i} = \dfrac{(1+i)^n - 1}{i}$\\[1.5pt]
\textbf{Konversi Frekuensi Berbeda} (kupon $m$/thn, yield $i^{(p)}$):
\begin{align*}
    \mathbf{i_{\text{kupon}}} &= \left(1 + \frac{i^{(p)}}{p}\right)^{p/m} - 1 \quad\Big(\text{sem vs } i^{(4)} \Rightarrow \big(1+\tfrac{i^{(4)}}{4}\big)^2 - 1\Big)
\end{align*}
\end{cbox}

% KOTAK 2
\begin{cbox}{2. Empat Rumus Harga Obligasi \& Status Premium/Diskon}
\begin{enumerate}[leftmargin=3.5mm, topsep=0.5pt, itemsep=1.2pt]
    \item \textbf{Dasar:} $\mathbf{P = Fr\,a_{\angl{n}i} + C v^n = Fr\,a_{\angl{n}i} + K}$
    \item \textbf{Prem/Diskon:} $\mathbf{P = C + (Fr - Ci)a_{\angl{n}i} = C + C(g - i)a_{\angl{n}i}}$
    \item \textbf{Base Amount:} $\mathbf{P = G + (C - G)v^n}$ \quad (paling cepat di kalkulator!)
    \item \textbf{Makeham:} $\mathbf{P = K + \dfrac{g}{i}(C - K)}$ \quad (pakai jika $n$ tak diketahui / serial)
\end{enumerate}
\textbf{Zero-Coupon:} $P = Cv^n \Rightarrow i = (C/P)^{1/n} - 1$. \quad \textbf{Perpetual:} $P = \dfrac{Fr}{i}$.\\[1.2pt]
\textbf{Kriteria Status (Cek $g$ vs $i$, BUKAN $r$ vs $i$ saat $C \neq F$!):}\\
$\bullet$ \textbf{Premium ($P > C$):} $g > i$ ($Fr > Ci$), \quad Prem $= P - C = C(g - i)a_{\angl{n}i}$\\
$\bullet$ \textbf{Diskon ($P < C$):} $g < i$ ($Fr < Ci$), \quad Disk $= C - P = C(i - g)a_{\angl{n}i}$
\end{cbox}

% KOTAK 3
\begin{gbox}{3. Jadwal Amortisasi Obligasi \& Rumus Langsung}
Nilai buku bergerak dari $B_0 = P$ menuju $B_n = C$.\\[1pt]
\textbf{Relasi Rekursif (isi tabel periode ke-$t$):}
\begin{align*}
    I_t &= i B_{t-1}, \qquad P_t = Fr - I_t, \\
    B_t &= B_{t-1} - P_t, \qquad B_{t+1} = B_t(1+i) - Fr
\end{align*}
$\bullet$ \textbf{Premium ($g > i$):} $P_t > 0 \Rightarrow B_t$ turun (\textit{Writing Down}).\\
$\bullet$ \textbf{Diskon ($g < i$):} $P_t < 0 \Rightarrow B_t = B_{t-1} + |P_t|$ naik (\textit{Writing Up}).\\[1.2pt]
\textbf{Rumus Langsung ke-$t$ (tanpa buat tabel dari awal!):}
\begin{align*}
    \mathbf{B_t} &= C + (Fr - Ci)a_{\angl{n-t}i}, \qquad \mathbf{P_t} = (Fr - Ci)v^{n-t+1} \\
    \mathbf{I_t} &= Fr - P_t = Fr - (Fr - Ci)v^{n-t+1}
\end{align*}
\textbf{Sifat Barisan Geometri Kolom Pokok ($P_t$):}\\
$\bullet$ $\mathbf{P_t = P_k(1+i)^{t-k}}$ \quad(cari $i$ dari dua pokok, misal $P_7/P_4 = (1+i)^3$)\\
$\bullet$ $\sum_{t=1}^n P_t = P - C = P_1 s_{\angl{n}i}$, \qquad $\sum_{t=1}^n I_t = nFr - (P - C)$\\
$\bullet$ Selisih book value: $\mathbf{B_{a-1} - B_b = \sum_{t=a}^b P_t = P_a \, s_{\angl{b-a+1}i}}$
\end{gbox}

% KOTAK 4
\begin{rbox}{4. Valuasi Antar 2 Tanggal Kupon (PALING SERING KELUAR!)}
Transaksi di waktu $t + k$ dengan $0 < k < 1$. Selalu berlaku:
\[
    \mathbf{B_{t+k}^f = B_{t+k}^m + Fr_k} \quad (\text{Flat/Dirty = Market/Clean + Accrued})
\]
\begin{center}
\scriptsize
\renewcommand{\arraystretch}{1.28}
\begin{tabular}{@{}p{1.75cm}@{\hspace{1.8mm}}p{2.05cm}@{\hspace{1.8mm}}p{1.85cm}@{\hspace{1.8mm}}p{1.95cm}@{}}
\toprule
\textbf{Komponen} & \textbf{Teoritis} & \textbf{Praktis} & \textbf{Semi-teoritis} \\
\midrule
Flat ($B_{t+k}^f$) & $B_t(1+i)^k$ & $B_t(1+ik)$ & $B_t(1+i)^k$ \\
Accrued ($Fr_k$) & $Fr\frac{(1+i)^k - 1}{i}$ & $k \cdot Fr$ & $k \cdot Fr$ \\
Market ($B_{t+k}^m$) & $B_{t+k}^f - Fr_k$ & $(1-k)B_t + kB_{t+1}$ & $B_t(1+i)^k - kFr$ \\
\bottomrule
\end{tabular}
\end{center}
$\bullet$ \textbf{Konvensi $k$:} \textit{30/360} (1 bln = 30 hr, 1 sem = 180 hr) atau \textit{Actual/Actual}.\\
$\bullet$ \textbf{Cek Premium/Diskon:} Gunakan $\mathbf{B_{t+k}^m - C}$ (BUKAN $B_{t+k}^f - C$!).
\end{rbox}

% KOTAK 5
\begin{abox}{5. Callable \& Putable Bond (Termasuk Call Bertingkat)}
\begin{center}
\scriptsize
\renewcommand{\arraystretch}{1.18}
\begin{tabular}{@{}p{1.8cm}@{\hspace{2mm}}p{2.95cm}@{\hspace{2mm}}p{2.95cm}@{}}
\toprule
\textbf{Kriteria} & \textbf{Callable ($\min_n P_n$)} & \textbf{Putable ($\max_n P_n$)} \\
\midrule
Pemegang Opsi & Penerbit (skenario terburuk) & Investor (skenario terbaik) \\
$C$ tetap, $g > i$ & \textbf{Call terawal} (pasca \textit{lockout}) & \textbf{Maturity} (paling akhir) \\
$C$ tetap, $g < i$ & \textbf{Maturity} (paling akhir) & \textbf{Put terawal} (pasca \textit{lockout}) \\
\bottomrule
\end{tabular}
\end{center}
\textbf{Resep Call Premium Bertingkat} (saat $C_n$ berubah antar segmen):\\
1. Di tiap segmen hitung $g_{\text{seg}} = Fr / C_{\text{seg}}$.\\
2. Jika $g_{\text{seg}} > i$ (prem) $\to$ ambil $n$ \textbf{terkecil} segmen; jika $g_{\text{seg}} < i$ (disk) $\to$ ambil $n$ \textbf{terbesar} segmen.\\
3. Hitung $P_n$ tiap kandidat (+ maturity) $\to$ ambil \textbf{nilai terkecil} ($\min$).\\[1pt]
\textbf{Soal Invers \& Cek Realisasi Yield:} Jika harga beli $P$ menjamin yield min $i^*$ di $n^*$, lalu realisasi di-\textit{call} pada $n_{\text{real}}$, bandingkan $P_{n_{\text{real}}}(i^*)$ vs $P$: jika $P_{n_{\text{real}}}(i^*) > P \Rightarrow i_{\text{aktual}} > i^*$.
\end{abox}

\clearpage

% ==============================================================================
% HALAMAN 2: SISI BELAKANG (TIER 2 & 3 — VARIASI, RUMUS LANJUTAN & JEBAKAN)
% ==============================================================================
\twocolumn[{
\begin{tcolorbox}[
    enhanced,
    colback=primaryblue,
    colframe=primaryblue,
    arc=1.2mm,
    boxrule=0pt,
    left=3mm, right=3mm, top=1.6mm, bottom=1.6mm,
    after skip=2mm
]
\begin{center}
    {\color{white}\normalsize\textbf{\textsf{CHEATSHEET 1 LEMBAR BOLAK-BALIK --- SISI BELAKANG (TIER 2 \& 3: VARIASI \& JEBAKAN)}}}\\[0.15em]
    {\color{amberbg}\scriptsize\textsf{Aproksimasi YTM $\bullet$ Reinvestasi $\bullet$ Kupon/Yield Tidak Konstan $\bullet$ Serial Bonds $\bullet$ Garis Lurus \& SF $\bullet$ Jebakan Maut}}
\end{center}
\end{tcolorbox}
}]

% KOTAK 6
\begin{cbox}{6. Aproksimasi YTM \& Dampak Perubahan Yield}
Diketahui harga $P$, nilai tebus $C$, tenor $n$ periode kupon, $g = Fr/C$:
\[
    \mathbf{k = \frac{P - C}{C}}, \qquad \mathbf{i \approx \frac{g - \dfrac{k}{n}}{1 + \dfrac{n+1}{2n}k}}, \qquad \text{(jika } n \text{ besar: } i \approx \frac{g - k/n}{1 + k/2}\text{)}
\]
$\bullet$ Hasil $i$ adalah yield \textbf{per periode kupon}; nominal tahunan $\mathbf{i^{(m)} = m \cdot i}$.\\
$\bullet$ \textbf{Dampak perubahan yield:} Hitung $P_{\text{baru}}$ di $i_{\text{baru}}$, lalu $\%\Delta P = \frac{P_{\text{baru}} - P}{P} \times 100\%$.
\end{cbox}

% KOTAK 7
\begin{cbox}{7. Reinvestasi Kupon (Overall Effective Yield)}
Setiap kupon $Fr$ ($n$ kali) ditabung pada suku bunga $j_{\text{tab}}$ per periode kupon hingga jatuh tempo (akhir tahun ke-$T$). Saldo akhir:
\[
    \mathbf{FV = Fr \, s_{\angl{n}j_{\text{tab}}} + C}, \qquad \mathbf{y_{\text{eff tahunan}} = \left(\frac{FV}{P}\right)^{1/T} - 1}
\]
\end{cbox}

% KOTAK 8
\begin{gbox}{8. Kupon Tidak Konstan \& Yield Rate Tidak Konstan}
\textbf{A. Kupon Bertahap (\textit{Step-Up / Step-Down}):} Pecah per blok anuitas:\\
Kupon $Q_1$ ($n_1$ per.), $Q_2$ ($n_2$ per.), $Q_3$ ($n_3$ per.), total $N = n_1+n_2+n_3$:
\[
    \mathbf{P_0 = Q_1 a_{\angl{n_1}i} + Q_2 v^{n_1} a_{\angl{n_2}i} + Q_3 v^{n_1+n_2} a_{\angl{n_3}i} + C v^N}
\]
(\textit{Book value} $B_t$ di tengah = PV sisa kupon ke depan + PV dari $C$).\\[1.5pt]
\textbf{B. Kupon Geometrik (\textit{Inflation-Adjusted}):} Kupon ke-$t$ = $Fr_1(1+k)^{t-1}$:
\begin{align*}
    \mathbf{P} &= Fr_1 \left[\frac{1 - \left(\tfrac{1+k}{1+i}\right)^n}{i - k}\right] + C(1+i)^{-n} \quad (k \neq i) \\
    \text{jika } k &= i:\quad \mathbf{P = n\,Fr_1 v + C v^n}
\end{align*}
\textbf{C. Yield Berubah ($i_1$ selama $n_1$ per., lalu $i_2$ selama $n_2$ per.):}
\[
    \mathbf{P = Fr \left[a_{\angl{n_1}i_1} + (1+i_1)^{-n_1} a_{\angl{n_2}i_2}\right] + C(1+i_1)^{-n_1}(1+i_2)^{-n_2}}
\]
\end{gbox}

% KOTAK 9
\begin{gbox}{9. Serial Bonds dengan Rumus Makeham (Slide 14 Hal. 10--12)}
Untuk $m$ obligasi berseri dengan $g = Fr_j/C_j$ seragam dan jatuh tempo $n_1,\dots,n_m$:
\[
    \mathbf{P' = K' + \frac{g}{i}(C' - K')}, \qquad C' = \sum_{j=1}^m C_j, \qquad K' = \sum_{j=1}^m C_j v^{n_j}
\]
\textbf{Trik Cepat Penebusan Merata} ($m$ angsuran $C_{\text{unit}}$ pada akhir tahun $a+1$ s.d. $a+m$):
\[
    \mathbf{C' = m \cdot C_{\text{unit}}}, \qquad \mathbf{K' = C_{\text{unit}} \, v^a \, a_{\angl{m}i}}
\]
\end{gbox}

% KOTAK 10
\begin{cbox}{10. Metode Garis Lurus \& Sinking Fund Obligasi ($j \neq i$)}
$\bullet$ \textbf{Garis Lurus (\textit{Straight Line}):} $P_t = \dfrac{P - C}{n}$ (konstan), $I_t = Fr - P_t$, $B_t = P - t P_t$.\\
$\bullet$ \textbf{Sinking Fund Obligasi ($j \neq i$):} Penyisihan premium ditabung di SF berbunga $j$ per periode kupon:
\[
    (Fr - iP_{\text{SF}})s_{\angl{n}j} = P_{\text{SF}} - C \implies \mathbf{P_{\text{SF}} = \frac{C\left(1 + g\,s_{\angl{n}j}\right)}{1 + i\,s_{\angl{n}j}}}
\]
(Jika $j < i$, maka $P_{\text{SF}} < P_{\text{standar}}$; jika $j = i$, kembali ke rumus standar).
\end{cbox}

% KOTAK 11
\begin{abox}{11. Ringkasan Klasifikasi Teori Singkat (Slide 12 Hal. 3--6)}
$\bullet$ \textbf{Accumulation Bond:} $C = F(1+r_a)^n$;\quad \textbf{Coupon Bond:} bayar kupon berkala.\\
$\bullet$ \textbf{Registered Bond:} pemilik tercatat resmi;\quad \textbf{Bearer Bond:} tidak tercatat (pemegang fisik = pemilik legal).\\
$\bullet$ \textbf{Mortgage Bond:} dijamin agunan (\textit{collateral}), lebih aman saat gagal bayar;\quad \textbf{Debenture Bond:} dijamin kredit umum tanpa agunan khusus.
\end{abox}

% KOTAK 12
\begin{rbox}{12. Checklist 4 Langkah \& 7 Jebakan Maut (WAJIB CEK!)}
\textbf{Checklist Kerja:} (1) Samakan $n, r, i$ ke satuan kupon $\to$ (2) Hitung $g = Fr/C$ \& cek $g$ vs $i$ $\to$ (3) Pilih rumus/metode $\to$ (4) Cek satuan jawaban ($i$ vs $i^{(2)}$; market vs flat).
\begin{enumerate}[leftmargin=3.5mm, topsep=0.5pt, itemsep=0.8pt]
    \item \textbf{$C \neq F$:} Status pakai \textbf{$g$ vs $i$}, JANGAN $r$ vs $i$!
    \item \textbf{Beda Frekuensi:} Kupon $m=2$, yield $i^{(4)} \Rightarrow$ konversi dulu ke semester.
    \item \textbf{Antar Kupon:} Premium/diskon = $\mathbf{B_{t+k}^m - C}$, BUKAN $B_{t+k}^f - C$.
    \item \textbf{Tanda $P_t$ Diskon:} $P_t < 0 \Rightarrow B_t = B_{t-1} - P_t$ harus \textbf{NAIK} ke $C$.
    \item \textbf{Pangkat $P_t$:} Pangkatnya $\mathbf{-(n - t + 1)}$, bukan $-(n - t)$!
    \item \textbf{Callable Bertingkat:} Uji tiap segmen $C_{\text{seg}}$ pakai $g_{\text{seg}}$, jangan tebak ujung!
    \item \textbf{Dua Arti $k$:} Di Bab 13 awal $k \in (0,1)$ (fraksi hari); di YTM $k = \frac{P-C}{C}$.
\end{enumerate}
\end{rbox}

\end{document}
